\begin{afterword}
\summaryhead

The post asks the reader to suppose it were proved, beyond doubt, that nothing is moral or right: whatever the basis of their moral philosophy, it has been torn apart, nothing can replace it, and all utilities equal zero. The author says this is not a real claim, and asks the reader to imagine concretely what they would do.

A list of questions follows. Would you still tip cab drivers, stay faithful, drag a fainted child off the train tracks? Would you eat the cheapest food, or the most expensive? Would you wear black and write gloomy poetry, which would only be a cached thought? Would you get out of bed, keep reading the blog, keep trying to be rational? The post ends by asking the reader to close their eyes and answer: what would you do, if nothing were right?

\responsehead

The post is a thought experiment in the form of questions, and it withholds the answers. Most readers will find their answers unchanged: they would still pull the child off the tracks. The questions about food show that the supposition gives no guidance either way, and the question about gloomy poetry heads off the stock picture of the nihilist. The point, which the post leaves for the reader to draw, is that much of what people do does not rest on their theory of what morality is.

The author gave the answers the next day, in ``The Moral Void,'' which the book leaves out. There the author wrote that the question drew numerous objections, ``as well there should have been.'' The next post in the book, ``Changing Your Metaethics,'' gives the conclusion in one line: if your metaethic stops telling you to save lives, you can drag the child off the tracks anyway.

The supposition joins two claims that philosophers keep apart. Moral error theorists, since Mackie in 1977, hold that nothing is morally right; they do not hold that everything is permissible, since that too would be a moral verdict. This does not change the thought experiment, which asks what the reader would do, not what would be allowed. Error theorists discuss a related question about moral language: whether to keep using it, as a kind of fiction, or to drop it.

In short: a short thought experiment in questions, close to one that moral error theorists also discuss, whose answers the author gave in a post the book leaves out and which the next post in the book summarizes.
\end{afterword}
